\id{ҒТАРМ 06.35.31}{}
\vspace{1em}
\begin{header}
\swa{Қ.Қ.~Мамутова, Ж.С.~Машраб, А.Н.~Айтымбетова}{ЖАСЫЛ ЭКОНОМИКА: ESG ҚАҒИДАТТАРЫНА СҮЙЕНГЕН ЖАСЫЛ ОБЛИГАЦИЯЛАР, ЖАСЫЛ НЕСИЕЛЕР ЖӘНЕ ИНВЕСТИЦИЯЛАР}

Қ.Қ.~Мамутова,
Ж.С.~Машраб\envelope,
А.Н.~Айтымбетова
\end{header}

\begin{affil}
М.Әуезов атындағы Оңтүстік Қазақстан университеті, Шымкент, Қазақстан

\corrauthor{Корреспондент-автор: zmashrab@bk.ru}
\end{affil}

Мақалада Қазақстанның жасыл экономикаға көшуі аясында ESG қағидаттарына
негізделген жасыл қаржыландыру құралдарының (жасыл облигациялар, жасыл
несиелер, әлеуметтік және ESG-облигациялар, жасыл инвестициялық қорлар,
жасыл микрокредиттер) рөлі мен даму келешегі қарастырылады. Зерттеудің
мақсаты-Қазақстан жағдайында жасыл облигациялар мен жасыл несиелердің
теориялық-әдіснамалық негіздерін айқындау, олардың жасыл экономиканы
және тұрақты дамуды қамтамасыз етудегі ықпалын талдау және осы
құралдарды дамыту бойынша институционалдық ұсыныстар ұсыну. Материалдық
база ретінде ғылыми әдебиеттер, ұлттық нормативтік-құқықтық және
стратегиялық құжаттар, халықаралық және қазақстандық статистикалық
деректер, эмитенттердің ашық ESG-есептілігі пайдаланылды. Әдіснамалық
тұрғыдан жүйелік және институционалдық тәсілдер, контент-талдау,
салыстырмалы және статистикалық талдау, құрылымдық-логикалық модельдеу
әдістері қолданылды.

Зерттеу нәтижесінде ESG өлшемдері мен қаржылық табиғатын ескеретін жасыл
қаржылық құралдардың матрицалық жіктемесі ұсынылды; әлемдік және
қазақстандық жасыл/ESG-облигациялар нарығының динамикасы сипатталып,
жасыл құралдардың корпоративтік облигациялар нарығындағы үлесінің әлі де
төмен екені, бірақ өсім қарқынының жоғары екені анықталды. Жаңартылатын
энергетика, энергия тиімділігі, коммуналдық инфрақұрылым, көлік, су
ресурстары және агросектор жасыл қаржыландырудың басым бағыттары ретінде
айқындалып, МЖӘ жобаларының қаржыландыру құрылымында жасыл компоненттің
әлеуеті көрсетілді. Верификация, ESG-есептілік стандарттары және
гринвошинг тәуекелдерімен байланысты институционалдық шектеулер
айқындалды. Авторлық қорытындылар жасыл қаржыландыруды ұлттық тұрақты
даму стратегияларымен, декарбонизация мақсаттарымен және жауапты
инвестиция мәдениетімен терең интеграциялау қажеттігін негіздейді және
болашақта эконометрикалық және эмпириялық зерттеулерге әдіснамалық база
бола алады.

{\bfseries Түйін сөздер:} эмитент, тұрақты даму, әлеуметтік облигациялар,
гринвошинг, жасыл өнімдер, жасыл жобалар, жасыл қаржыландыру қоры
\vspace{1em}
\begin{header}
ЗЕЛЕНАЯ ЭКОНОМИКА: ЗЕЛЕНЫЕ ОБЛИГАЦИИ, ЗЕЛЕНЫЕ КРЕДИТЫ И ИНВЕСТИЦИИ НА ОСНОВЕ ПРИНЦИПОВ ESG

К.К.~Мамутова,
Ж.С.~Машраб\envelope,
А.Н.~Айтымбетова
\end{header}

\begin{affil}
Южно-Казахстанский университет им. М. Ауэзова, Шымкент, Казахстан,

е-mail: zmashrab@bk.ru
\end{affil}

В статье рассматривается роль и перспективы развития инструментов
«зеленого» финансирования, основанных на принципах ESG (зеленые
облигации, зеленые кредиты, социальные и ESG-облигации, зеленые
инвестиционные фонды, зеленые микрокредиты) в контексте перехода
Казахстана к зеленой экономике. Цель исследования - выявление
теоретических и методологических основ зеленых облигаций и зеленых
кредитов в контексте Казахстана, анализ их влияния на обеспечение
зеленой экономики и устойчивого развития, а также предоставление
институциональных рекомендаций по развитию этих инструментов.
Материальная база включала научную литературу, национальные
нормативно-правовые и стратегические документы, международные и
казахстанские статистические данные, а также открытую ESG-отчетность
эмитентов. С методологической точки зрения использовались
системно-институциональный подход, контент-анализ,
сравнительно-статистический анализ и методы структурно-логического
моделирования.

В исследовании представлена матричная классификация «зеленых» финансовых
инструментов, учитывающая ESG-факторы и финансовую природу; описана
динамика мирового и казахстанского рынков «зеленых»/ESG-облигаций, и
установлено, что доля «зеленых» инструментов на рынке корпоративных
облигаций остается низкой, но темпы роста высоки. В качестве
приоритетных областей «зеленого» финансирования были определены
возобновляемая энергия, энергоэффективность, муниципальная
инфраструктура, транспорт, водные ресурсы и агропромышленный сектор, а
также продемонстрирован потенциал «зеленой» составляющей в структуре
финансирования проектов ГЧП. Выявлены институциональные ограничения,
связанные с верификацией, стандартами отчетности по ESG и рисками
«зеленого отмывания». Выводы автора обосновывают необходимость глубокой
интеграции «зеленого» финансирования в национальные стратегии
устойчивого развития, цели декарбонизации и культуру ответственного
инвестирования, и могут служить методологической основой для будущих
эконометрических и эмпирических исследований.

{\bfseries Ключевые слова:} эмитент, устойчивое развитие, социальные
облигации, гринвошинг, экологически чистая продукция, экологические
проекты, фонд зеленого финансирования
\vspace{1em}
\begin{header}
GREEN ECONOMY: GREEN BONDS, GREEN LOANS, AND ESG-BASED INVESTMENTS

K.K.~Mamutova,
Zh.S.~Mashrab\envelope,
A.N.~Aitymbetova
\end{header}

\begin{affil}
M. Auezov South Kazakhstan University, Shymkent, Kazakhstan,

е-mail: zmashrab@bk.ru
\end{affil}

This article examines the role and development prospects of green
financing instruments based on ESG principles (green bonds, green loans,
social and ESG bonds, green investment funds, and green microloans) in
the context of Kazakhstan' s transition to a green
economy. The aim of the study was to identify the theoretical and
methodological foundations of green bonds and green loans in the context
of Kazakhstan, analyze their impact on ensuring a green economy and
sustainable development, and provide institutional recommendations for
the development of these instruments. The material base included
scientific literature, national regulatory and strategic documents,
international and Kazakhstani statistical data, and
issuers'{} public ESG reporting. From a methodological
perspective, a systems-institutional app\-roach, content analysis,
comparative statistical analysis, and structural-logical modeling
methods were used.

The study presents a matrix classification of green financial
instruments, taking into account ESG factors and their financial nature.
It describes the dynamics of the global and Kazakhstani green/ESG bond
markets, and establishes that the share of green instruments in the
corporate bond market remains low, but growth rates are high. Renewable
energy, energy efficiency, municipal infrastructure, transport, water
resources, and the agro-industrial sector were identified as priority
areas for green financing, and the potential of the green component in
the structure of PPP project financing was demonstrated. Institutional
constraints related to verification, ESG reporting standards, and
greenwashing risks were identified. The author' s
conclusions substantiate the need for deep integration of green finance
into national sustainable development strategies, decarbonization goals,
and a culture of responsible investment, and can serve as a
methodological basis for future econometric and empirical research.

{\bfseries Keywords:} issuer, sustainable development, social bonds,
greenwashing, environmentally friendly products, environmental projects,
green finance fund

\begin{multicols}{2}
{\bfseries Кіріспе.} ХХІ ғасыр жағдайында климаттың өзгеруі, табиғи
ресурстар қорының азаюы және жаһандық экологиялық тәуекелдердің күшеюі
әлемдік шаруашылықты дәстүрлі индустриялық үлгіден жасыл экономикаға бет
бұруға мәжбүрлеуде. Жасыл экономика экономикалық өсуді қоршаған ортаға
түсетін жүктемені төмендетумен, ресурстарды ұқыпты пайдаланумен және
әлеуметтік теңдік қағидаттарымен ұштастыруға бағытталған. Осындай
жағдайда қаржы жүйесінің жаңаруы -- жасыл облигациялар, жасыл несиелер,
сондай-ақ әлеуметтік және ESG облигациялары сияқты қаржы құралдарының
кеңеюі- тұрақты даму мақсаттарына жетудің маңызды тетіктерінің біріне
айналып отыр {[}1{]}.

Қазақстанда да жасыл қаржыландыру нарығы қалыптасу кезеңінен өтіп, АХҚО,
Astana International Exchange (AIX), AIFC Green Finance Centre, «Даму»
қоры, Қазақстан Даму Банкі сияқты институттардың қатысуымен жасыл
облигациялар мен жасыл несиелерді орналастыру тәжірибесі біртіндеп
кеңейіп келеді. AIX биржалық алаңында алғашқы жасыл облигациялар 2020
жылдан бастап тіркеліп, жаңартылатын энергия көздері, энергия тиімділігі
және өзге де экологиялық сипаттағы бастамаларды қаржыландыруға
бағытталды {[}2, 3{]}. Мұндай бағыт Қазақстанның төмен көміртекті даму
жолын таңдауы және ESG қағидаттарын кезең-кезеңімен ұлттық қаржы
жүйесіне енгізу ниетімен үйлеседі.

Сонымен бірге жасыл қаржылық өнімдерге деген сұраныстың өсуі бірқатар
ішкі қайшылықтарды да айқын көрсетіп отыр. Бір жағынан, жасыл
облигациялар мен несиелер эмитенттерге капиталдың қосымша көздерін ашып,
инвесторларға тәуекелі мен кірістілігі салыстырмалы түрде теңгерілген
ұзақ мерзімді жобаларға инвестиция салуға мүмкіндік береді. Екінші
жағынан, нарықта гринвошинг қатері, жобалардың «жасыл» мәртебесін
айқындайтын өлшемдердің бірізді еместігі, верификация рәсімдері мен
ESG-есептіліктің жеткілікті деңгейде қалыптаспауы байқалады. Осыған
байланысты жасыл облигациялар, жасыл несиелер және ESG-инвестициялар
саласындағы ұлттық және халықаралық тәжірибені ғылыми негізде жүйелеу,
олардың Қазақстанның тұрақты дамуына қосатын үлесін талдау және
институционалдық шектеулерді анықтау ерекше маңызға ие болып отыр.

Осыған орай зерттеу тақырыбын айқындаудың негізі жасыл экономикаға көшу
жағдайында қаржы нарығында пайда болған «жасыл» және ESG-бағдарланған
жаңа құралдардың рөлі теориялық және эмпирикалық тұрғыдан әлі жеткілікті
ашылмағанымен, сондай-ақ Қазақстанда жасыл облигациялар мен жасыл
несиелер нарығының әлі де қалыптасу сатысында, салыстырмалы түрде жаңа
құбылыс болуымен түсіндіріледі. Зерттеу тақырыбының өзектілігі жасыл
инвестициялар есебінен ұлттық экономиканы құрылымдық жаңғыртуға,
көміртек ізі көлемін азайтуға, әлеуметтік маңызы бар бастамаларды
қаржыландыруға және халықаралық тұрақты даму күн тәртібімен (SDGs)
үйлесуге нақты мүмкіндіктер ашатынынан көрінеді.

Зерттеудің мақсаты-Қазақстан жағдайында ESG қағидаттарына сүйенетін
жасыл облигациялар, жасыл несиелер мен инвестициялардың
теориялық-әдіснамалық негіздерін нақтылау, олардың жасыл экономиканы
қалыптастыру және тұрақты дамуды қолдаудағы маңызын зерделеу, сондай-ақ
аталған құралдарды одан әрі өркендетуге бағытталған ұсынымдар әзірлеу.

Осы мақсатқа сәйкес зерттеуде мынадай міндеттер шешіледі:

- жасыл экономика, жасыл қаржыландыру, ESG-облигациялар және әлеуметтік
облигациялар бойынша ғылыми көзқарастарды жүйелеу және ұғымдық-санаттық
аппаратты нақтылау;

- ICMA Green Bond Principles, Climate Bonds Standard және өзге де
құжаттар аясындағы жасыл облигациялар мен жасыл несиелердің халықаралық
стандарттарын Қазақстандағы қалыптасқан тәжірибемен салыстырмалы талдау;

- Қазақстанда жасыл облигациялар, жасыл несиелер, сондай-ақ әлеуметтік
және ESG-облигациялар нарығының қалыптасу ерекшеліктерін, басты
эмитенттерін және қаржыландырылатын жасыл жобалардың бағыттарын
сипаттау;

- жасыл қаржылық өнімдерге қатысты гринвошинг қатерлерін, сыртқы
верификация рәсімдерін және ESG-есептілік тетіктерін айқындау;

- жасыл облигациялар мен несиелердің эмитенттер, инвесторлар және қоғам
үшін экономикалық, экологиялық әрі әлеуметтік нәтижелілігін бағалау;

- жасыл қаржыландыру құралдарын одан әрі жетілдіру және оларды ұлттық
тұрақты даму стратегияларымен үйлестіру бойынша институционалдық және
реттеушілік ұсыныстар әзірлеу.

Зерттеудің әдістемелік негізін жүйелік, институционалдық және
салыстырмалы талдау құрайды; ғылыми әдебиеттерге шолу,
нормативтік-құқықтық және стратегиялық құжаттарды контент-талдау,
халықаралық және ұлттық статистика мен эмитенттердің ашық ESG-есептілігі
және жасыл облигациялар шығарылымы туралы деректерді талдау әдістері
пайдаланылады. Зерттеудің гипотезасы: егер Қазақстанда ESG қағидаттарына
сүйенген жасыл облигациялар мен жасыл несиелер нарығының
институционалдық және реттеушілік базасы нығайып, верификация мен
ашықтық тетіктері дамыса, бұл жасыл экономикаға көшу қарқынын
жеделдетіп, ұлттық тұрақты даму мақсаттарына жетуге және қаржы жүйесінің
орнықтылығын арттыруға ықпал етеді. Теориялық маңыздылығы -- жасыл
экономика, жасыл қаржыландыру және ESG-инвестициялар жөніндегі
түсініктерді нақтылау және жасыл облигациялар мен несиелерді әлеуметтік
және ESG облигацияларымен өзара байланыста қарастыруда; практикалық
маңыздылығы -- алынған қорытындыларды мемлекеттік органдар мен қаржы
институттарының жасыл қаржыландыру саясатын жетілдіруде, жасыл қаржы
қорларының қызметін ұйымдастыруда және ЖОО-лардың оқу жоспарларына жасыл
қаржы мен ESG бойынша пәндер енгізуде қолдану мүмкіндігінде.

{\bfseries Материалдар мен әдістер.} Зерттеудің әдіснамалық негізі негізгі
сұрақ төңірегінде құрылған: ESG-ге бағытталған жасыл қаржы құралдары
(жасыл облигациялар, жасыл несиелер, әлеуметтік және ESG облигациялары)
Қазақстанның жасыл экономикаға көшуіне қалай ықпал ете алады және бұл
үшін қандай институционалдық алғышарттар қажет. Негізгі гипотеза
келесідей тұжырымдалған: егер Қазақстандағы жасыл облигациялар мен несие
нарығының институционалдық және нормативтік базасы нығайтылса, ашықтық,
тиісті тексеру және стандартталған ESG есептілігі қамтамасыз етілсе, бұл
құралдар төмен көміртекті даму моделіне қарай ілгерілеуді айтарлықтай
жеделдетіп, қаржы жүйесінің тұрақтылығын арттыра алады.

Зерттеудің эмпирикалық базасына мыналар кіреді:

- 2015-2025 жылдарға арналған жасыл экономика, жасыл қаржы және ESG
инвестициялары бойынша қазақ, орыс және ағылшын тілдеріндегі ғылыми
еңбектер (монографиялар, рецензияланған журналдардағы мақалалар,
аналитикалық шолулар);

- тұрақты дамуды, жасыл экономикаға көшуді, бағалы қағаздар нарығының
жұмыс істеуін және ESG ақпаратын ашуды реттейтін Қазақстан
Республикасының нормативтік актілері мен стратегиялық құжаттары;

- әлемдік және Қазақстандағы жасыл және ESG облигацияларын орналастыру
көлемі, корпоративтік облигациялар сегментінің құрылымы және PPP/MZhE
жобаларының қаржылық параметрлері туралы статистика (үкіметтік және
биржа дерекқорларынан алынған деректер, AIFC, AIX, ұлттық реттеушілер
мен халықаралық институттардың есептері);

- жасыл облигациялар эмитенттері мен несиелік мекемелердің жария ESG
есеп беруі. Пайдаланылған материал сапалық компонентті (құралдардың
түрлерін, реттеуші режимдерді, институционалдық қатысушыларды сипаттау)
және сандық компонентті (орналастыру көлемдері, нарық үлестері және
«жасыл» жобалардың секторлары бойынша динамикалық қатарлар мен
құрылымдық сипаттамалар) біріктіреді.

Зерттеу бірнеше өзара байланысты кезеңдерден өтті:

- жасыл экономикаға, жасыл қаржыға және ESG инвестициялауға негізгі
тәсілдерге теориялық және тұжырымдамалық шолу жүргізілді, тұжырымдамалық
негіз нақтыланды;

- тұрақты даму және қаржы нарығы саласындағы Қазақстанның ұлттық
нормативтік базасы мен стратегиялық құжаттарына институционалдық талдау
жүргізілді;

- ESG өлшемдері мен қаржылық сипаты (қарыз, меншікті капитал, гибридті,
бюджеттік/субсидияланған) бойынша жасыл қаржы құралдарының меншікті
матрицалық жіктемесі әзірленді;

- әлемдік және қазақстандық жасыл және ESG облигациялары нарықтары
туралы ақпарат жиналып, статистикалық өңделді, динамикалық қатарлар мен
салыстырмалы көрсеткіштер (корпоративтік облигация сегментіндегі жасыл
құралдардың үлесі және MZHӘ жобаларын қаржыландыру көздерінің құрылымы)
есептелді;

- нәтижелер халықаралық зерттеулермен салыстырылды және Қазақстанда
жасыл қаржыландыруды одан әрі дамыту бойынша ұсыныстар тұжырымдалды.

Зерттеу жүйелік және институционалдық тәсілге негізделген, бұл жасыл
қаржыны тұрақты даму мен қаржылық реттеудің кеңірек жүйесінің ажырамас
бөлігі ретінде қарастыруға мүмкіндік береді. Келесі әдістер қолданылды:

- негізгі үрдістер мен нормативтік талаптарды анықтау үшін ғылыми
дереккөздердің, нормативтік құжаттардың және стратегиялық
бағдарламалардың мазмұндық талдауы;

- жасыл және ESG облигациялары нарықтарының жұмыс істеуіндегі
халықаралық және ұлттық тәжірибелердің салыстырмалы талдауы;

- стандартты кеңсе бағдарламалық жасақтамасынан кестелік және
визуализация құралдарын пайдалана отырып, уақыт қатарларын статистикалық
өңдеу (өсу қарқынын, құрылымдық үлестерді бағалау және деректерді
графиктер мен диаграммаларда ұсыну);

- MZhÄ жобаларын қаржыландыру көздерінің схемасын және жасыл қаржы
құралдарының матрицалық жіктемесін әзірлеудегі құрылымдық және логикалық
модельдеу;

- анықталған үрдістер мен институционалдық кедергілерді түсіндіруде
қолданылатын сараптамалық бағалау элементтері.

Қолданылған материалдар мен әдістердің үйлесімі бізге мыналарға
мүмкіндік берді: жасыл қаржы құралдарының матрицалық жіктелуінің
авторлық нұсқасын негіздеуге; Қазақстанның жаһандық жасыл және ESG
облигациялар нарығындағы орнын, сондай-ақ корпоративтік облигациялар
сегментіндегі «жасыл» құралдардың үлесін сандық бағалауға; «жасыл»
жобаларды іске асырудың басым секторларын және оларды қаржыландыру
көздерінің конфигурациясын анықтауға; тексеруге, ESG есептілігіне және
«жасыл жуу» қаупіне байланысты негізгі институционалдық кемшіліктерді
анықтауға. Жалпы алғанда, бұл зерттеу әдіснамасын қойылған мақсаттарға
сәйкес келетіндей қарастыруға, алынған нәтижелердің сенімділігін
қамтамасыз етуге және Қазақстандағы жасыл қаржы және ESG инвестициялық
саясаты саласындағы кейінгі, егжей-тегжейлі эмпирикалық және
эконометрикалық зерттеулер үшін негіз қалыптастыруға мүмкіндік береді.

{\bfseries Нәтижелер мен талқылау.} Зерттеу нәтижесінде жасыл экономика мен
ESG қағидаттарына сүйенетін жасыл қаржыландыру құралдарының халықаралық
және ұлттық тәжірибесі жан-жақты қарастырылып, Қазақстан жағдайына
икемделген бірқатар қорытындылар жасалды.

Ең алдымен теориялық еңбектер мен нормативтік құжаттар негізінде жасыл
қаржы құралдарының (жасыл облигациялар, жасыл несиелер, әлеуметтік және
ESG-облигациялар, жасыл инвестициялық пай қорлары, жасыл микрокредиттер
және т.б.) құрылымдық жіктемесі жасалды {[}4,5{]}. Ұсынылған жіктеу екі
басты өлшемге сүйенеді: ESG құрамдастары (E-экологиялық ықпал,
S-әлеуметтік нәтиже, G-корпоративтік басқару деңгейі) және қаржылық
табиғаты (қарыздық, үлестік, аралас нысандар, сондай-ақ
бюджеттік/субсидиялық құралдар). Осы тәсілдің негізінде «Жасыл қаржылық
құралдардың ESG өлшемдері мен қаржылық табиғаты бойынша жіктелуі» атты
матрица түріндегі визуалды модель құрастырылды (сурет 1). Мұндай жүйелеу
экологиялық тұрғыдан қауіпсіз өнім желілерін қалыптастыруда эмитенттер
мен реттеуші органдарға әдістемелік бағдар ретінде қызмет ете алады.
\end{multicols}

\fig{e/image11}[1-сурет. Жасыл қаржылық құралдарды ESG өлшемдері мен қаржылық
сипаты негізінде матрицалық түрде топтастыру]

\begin{multicols}{2}
Сурет 1-де берілген матрицалық жіктеу жасыл қаржылық құралдардың негізгі
бөлігі қарыздық құралдарға тиесілі екенін, сондай-ақ басым назардың
экологиялық өлшемге (E) аударылатынын көрсетеді, ал әлеуметтік (S) және
корпоративтік басқару (G) қырлары, әсіресе дәстүрлі жасыл облигациялар
мен несиелерде, көбіне қосалқы немесе шектеулі түрде ескеріледі.
ESG-облигациялар мен жасыл инвестициялық қорлар үш өлшемді де
салыстырмалы түрде тең дәрежеде қамтитын интеграцияланған форматымен
ерекшеленіп, жауапты инвесторлар үшін тартымдылығы жоғары құралдар
қатарында көрінеді. Ұсынылған матрица эмитенттерге қай
ESG-компоненттерін күшейту қажеттігін айқындауға, ал реттеуші органдарға
ұлттық жасыл таксономияны әзірлеу барысында қандай құралдарды бірінші
кезекте экологиялық, қандайларын әлеуметтік немесе кешенді тұрақты даму
мақсаттарымен байланыстыру орынды екенін бағалауға мүмкіндік береді.

Халықаралық статистика мен Қазақстанға қатысты қолжетімді мәліметтерді
салыстырмалы талдау жасыл облигациялар нарығының тұрақты өсуін және
ESG-облигациялар сегментінің жедел кеңеюін айқындады {[}4,5{]}. Әлемдік
деңгейде жасыл және басқа да тұрақты облигациялар (social,
sustainability-linked) жиынтық нарығы соңғы жылдары дәстүрлі
облигациялардан жылдам өсіп, эмитенттер құрылымында инфрақұрылымдық
компаниялармен бірге қаржы ұйымдарының үлесі ұлғайғаны байқалады.
Қазақстанда 2020 жылдан кейін жасыл және ESG-облигациялар бойынша
алғашқы орналастырулар жүзеге асырылып, ұлттық биржалардағы олардың
көлемі жыл сайын артып келеді; аталған үрдіс Мангибаева Д.Д. және басқа
зерттеушілердің нәтижелерімен де қуатталады {[}5{]}. Осы қорытындыларды
визуализациялау үшін 2015-2024 жылдар аралығында әлемдік және
Қазақстандағы жасыл/ESG-облигациялар нарығы көлемінің салыстырмалы
динамикасы сурет 2-де берілген.

Суретте 2015-2024 жж. аралығындағы әлемдік және Қазақстандағы
жасыл/ESG-облигациялар нарығының динамикасы салыстырмалы түрде берілген:
жаһандық деңгейде бұл сегмент негізінен тұрақты өсу траекториясын
сақтап, 89 млрд доллардан 975 млрд долларға дейін ұлғайған, тек 2020
жылы ғана қысқа мерзімді төмендеу байқалады. Қазақстанда жасыл
облигациялар нарығы 2021 жылдан бастап қалыптасып, бастапқы өте төмен
көлемдерден (10 млн доллар) 2024 жылға қарай 475 млн долларға дейін
жедел өсім көрсеткен. Осылайша, абсолюттік көрсеткіштер тұрғысынан
ұлттық нарық әлемдік ауқымнан әлі де айтарлықтай төмен болғанымен, өсім
қарқынының жоғары болуы елде жасыл қаржыландыруды кеңейтуге елеулі
әлеует пен институционалдық алғышарттардың бар екенін білдіреді.
\end{multicols}

\fig[0.65\textwidth]{e/image20}[2-сурет.2015--2024 жылдары әлемдік және Қазақстандағы
жасыл/ESG-облигациялар нарығы көлемінің салыстырмалы өзгеру
динамикасы\\\normalfont{\emph{Ескертпе: суреттегі
жасыл/ESG-облигациялар жөніндегі әлемдік және Қазақстан көрсеткіштері
{[}4-5; 11{]} деректеріне негізделген авторлық талдау нәтижесінде
құрастырылған}}]

\begin{multicols}{2}
Сонымен бірге, Қазақстан бойынша алынған деректерді талдау жасыл
облигациялардың жалпы корпоративтік облигациялар нарығындағы үлесі әлі
де мардымсыз болғанымен, олардың өсу қарқыны жоғары екенін көрсетті. Бұл
үрдіс корпоративтік облигациялар нарығында жасыл және
ESG-облигациялардың пайыздық үлесін көрсететін диаграммада (сурет 3)
көрнекі түрде ұсынылған.
\end{multicols}

\fig[0.6\textwidth]{e/image21}[3-сурет. Қазақстандағы корпоративтік облигациялар
нарығындағы жасыл және ESG-облигациялардың үлесі,
\%\\\normalfont{\emph{Ескертпе: сурет Қазақстандағы
жасыл/ESG-облигациялар нарығы мен корпоративтік облигациялар
құрылымына қатысты деректер негізінде автор жүргізген есептеулерді
көрсетеді; есептеулерде {[}9-11{]} дереккөздері пайдаланылған.}}]

\begin{multicols}{2}
Сурет 3-тегі мәліметтер 2020-2024 жылдар аралығында Қазақстанның
корпоративтік облигациялар нарығындағы жасыл және ESG-облигациялар
үлесінің 0,1\%-дан 3,5\%-ға дейін өскенін, яғни үрдістің сапалық
тұрғыдан біртіндеп жоғарылағанын көрсетеді.2020-2022 жж. аралығында бұл
үлес небәрі 0,1-0,4\% шегінде баяу ұлғайып, жасыл сегменттің негізінен
пилоттық сипатта болғанын байқатса, 2023 жылдан бастап серпінді кеңею
кезеңі басталып, алдымен 1,4\%-ға дейін, ал алдын ала бағалау бойынша
2024 жылы 3,5\%-ға дейін өсуі жасыл қаржы құралдарының корпоративтік
қаржы құрылымына нақты ене бастағанын білдіреді. Сонымен қатар, олардың
үлесі жалпы нарық көлемінде әлі де салыстырмалы түрде төмен деңгейде
қалып отыр, бұл институционалдық қолдауды күшейту мен эмиссия ауқымын
ұлғайту қажеттігін көрсетеді.

Зерттеу қорытындылары Қазақстан жағдайында «жасыл» экономиканың басым
секторларын да айқындауға мүмкіндік берді. Жаңартылатын энергия көздері,
энергия тиімділігі, су ресурстарын басқару, қалдықтарды басқару, көлік
және аграрлық сектор жасыл қаржыландырудың негізгі векторлары ретінде
белгіленді {[}6,7{]}. Жагыпарова А.О. және оның әріптестерінің
жұмыстарын, сондай-ақ ұлттық стратегиялық құжаттарды талдау негізінде
жасыл қаржы құралдары есебінен қаржыландырылған немесе алдағы уақытта
қаржыландырылуы мүмкін жобалардың секторлық құрылымы жүйеленді {[}6{]}.
Бұл құрылымға жаңартылатын энергия және төмен көміртекті генерация,
тұрғын үй-коммуналдық шаруашылық (энергоүнемдеуші жаңғырту, «жасыл»
құрылыс), көлік саласы (электр көлігін дамыту, қоғамдық көлік паркін
жаңарту), су ресурстарын басқару және ауыл шаруашылығы (суды үнемдеу
технологиялары, топырақ құнарлылығын сақтау), сондай-ақ қалдықтарды
басқару мен қайта өңдеу бағыттары кіреді. Аталған секторлардың
салыстырмалы үлесі «Қазақстандағы жасыл қаржыландыру есебінен іске
асырылған (немесе жоспарланып отырған) жобалардың секторлық құрылымы»
атты диаграммада (сурет 4) көрнекі түрде көрсетілген. Бұл нәтижелер
ұлттық «жасыл» таксономияны әзірлеу үшін эмпирикалық база ретінде қызмет
етіп, экологиялық тұрғыдан таза жобалар мен өнімдерді іріктеу
өлшемшарттарын нақтылауға ықпал етеді.
\end{multicols}

\fig[0.6\textwidth]{e/image12}[4-сурет. Қаржыландыру көздері бойынша МЖӘ жобаларының құрылымдық
үлесі, \%\\\normalfont{\emph{Ескертпе: {[}6, 7, 11{]} дереккөздеріндегі мәліметтер негізінде авторлар әзірлеген}}]

\begin{multicols}{2}
Сурет 4 мәліметтері бойынша, МЖӘ жобаларының қаржыландыру құрылымында
негізгі үлесті қоғамдық бюджет (37\%) пен отандық жеке инвесторлар
(33\%) иеленеді, халықаралық даму институттары мен шетелдік
инвесторлардың үлесі 16\% шамасында. Жасыл қаржыландыру әзірге 14\%
деңгейінде қалып, жалпы портфельде қосалқы орынға ие болғанымен, оның
МЖӘ жобалары құрылымында жеке санат ретінде бөлінуі алдағы уақытта
«жасыл» құралдарды кеңейтіп, оларды қоғамдық және жеке қаржы көздерімен
тереңірек ұштастыруға елеулі мүмкіндік бар екеніне меңзейді.

Нормативтік және әдістемелік материалдарды (ұлттық «жасыл» экономикаға
көшу тұжырымдамасы, жасыл қаржыға арналған бағдарламалық актілер,
әдістемелік ұсынымдар) талдау Қазақстанда жасыл қаржыландыруды дамытуға
арналған институционалдық базаның кезең-кезеңімен күшейіп жатқанын
көрсетті {[}6{]}. Алайда жасыл облигациялар мен несиелерге қатысты
жобалардың «жасыл» мәртебесін растау (верификация) тәжірибесінің ICMA,
Climate Bonds Initiative және өзге де халықаралық стандарттарға толық
сәйкестендірілмеуі, ESG-есептіліктің бірыңғай форматының кеңінен
қолданылмауы, сондай-ақ «жасыл» және әлеуметтік әсер көрсеткіштерін
айқындау әдістемелерінің біріздендірілмеуі гринвошинг тәуекелдерін
күшейтетін факторлар ретінде анықталды {[}7{]}. Аталған тұжырымдар
5-бөлімде қарастырылған гринвошингтің негізгі көріністері мен оны шектеу
тетіктерін нақтылай түсіп, ұлттық деңгейде сыртқы верификаторлар
институтын, тәуелсіз ESG-рейтингтер жүйесін және ашықтыққа негізделген
есептілік стандарттарын дамыту қажеттігін дәлелдейді.

Алынған мәліметтер жасыл экономиканы қаржыландыруға арналған құралдар
мен институттарды жүйелеу, олардың нарықтық даму динамикасын сипаттау
және Қазақстанға тән ерекшеліктерін көрсету тұрғысынан бірқатар мәнді
тұжырымдар жасауға мүмкіндік береді. Ең алдымен, ESG өлшемдері мен
қаржылық сипаттарына сүйене отырып құрастырылған жасыл қаржы
құралдарының матрицалық жіктемесі халықаралық әдебиетте берілген
тәсілдермен әдістемелік тұрғыда үйлеседі. Шетелдік зерттеулерде жасыл
облигациялар мен өзге де тұрақты қаржы құралдары көбіне экологиялық,
әлеуметтік және басқарушылық әсерді қатар қамтитын кешенді өнімдер
ретінде сипатталатыны көрсетілген, бұл біздің құралдарды E, S, G
компоненттері мен қарыздық/үлестік сипатының қиылысында қарастыру
логикасымен сәйкес келеді {[}8{]}. Осындай жіктеу форматы эмитенттер
үшін өнімдер желісін жоспарлауда, ал реттеуші органдар үшін ұлттық жасыл
таксономияның құрылымын жобалауда қолданбалы бағдар қызметін атқара
алады.

Жасыл облигациялар мен ESG-құралдар нарығының жаһандық және ұлттық
деңгейдегі даму траекториясын салыстыру нәтижелері халықаралық әдебиетте
сипатталған негізгі бағыттармен үйлеседі. Ғылыми еңбектерде жасыл
облигациялар эмиссиясының өсуі эмитенттер үшін капитал құнының
төмендеуімен, ақпараттық асимметрияның азаюымен және нарықтық
өтімділіктің жақсаруыімен қатар жүретіні көрсетілген {[}8{]}. Біздің
зерттеу қорытындылары да әлемдік деңгейде тұрақты облигациялар нарығының
жалпы облигациялар нарығына қарағанда қарқындырақ кеңейіп жатқанын, ал
Қазақстанда жасыл және ESG-облигациялар сегменті көлемі жағынан әлі
шағын болғанымен, өсім қарқынының жоғары екенін аңғартады. Бұл, бір
жағынан, халықаралық үрдістердің ұлттық нарыққа белгілі бір уақыттық
кідіріспен трансляцияланатынын, екінші жағынан -- жасыл қаржыландыруды
ұлғайту үшін институционалдық «мүмкіндіктер терезесінің» сақталып
отырғанын білдіреді.

Сонымен қатар, ESG-инвестициялардың тәуекел мен табыстылық арақатынасы
жөніндегі халықаралық талқылаулар бұл құралдарға деген сұраныс артып
келе жатқанымен, инвесторлардың күтулері біркелкі емес екенін көрсетеді.
Жүйелік шолуларда ұзақ мерзімді кезеңде «жасыл» активтердің кірістілігі
дәстүрлі құралдарға қарағанда төмендеу болуы мүмкін екені, ал қысқа
мерзімде ESG-портфельдер белгілі бір артықшылықтар ұсына алатыны атап
өтіледі {[}9{]}. Біздің зерттеу қорытындылары да осы пайымдаулармен
үндес: Қазақстандағы эмиссияның құрылымы мен ауқымы әзірге шағын
болғандықтан, «жасыл премия» немесе «жасыл дисконт» құбылыстарын
статистикалық тұрғыдан нақтылы бағалау мүмкіндігі шектеулі, дегенмен
эмитенттер үшін беделдік тиімділік пен инвесторлар шеңберін кеңейту
факторы айқын байқалады.

Жасыл жобалардың секторлық бөлінісі де халықаралық тәжірибемен
салыстырғанда қисынды әрі негізделген екенін көрсетеді. Қазақстанда
жасыл қаржыландырудың негізгі бағыттарының жаңартылатын энергетика,
энергия тиімділігі, коммуналдық инфрақұрылым, көлік және агросекторға
шоғырлануы ұлттық «жасыл» экономикаға көшу тұжырымдамасымен, сондай-ақ
ESG қағидаттарын стратегиялық деңгейде кіріктіруге бағытталған зерттеу
нәтижелерімен үйлеседі {[}10{]}. Мұндай құрылым елдің
ресурстық-энергетикалық ерекшеліктеріне, көміртек қарқындылығының
деңгейіне және өңірлік даму басымдықтарына сай келеді, бірақ сонымен
бірге әлеуметтік әсері жоғары бастамалардың (қолжетімді «жасыл» тұрғын
үй, инклюзивті инфрақұрылым, ауылдық аумақтарды жаңғырту) үлесін
көбейтуге мүмкіндік беретін әлеует сақталып отыр.

Институционалдық ортаға қатысты нәтиже­лер біркелкі емес көріністі
байқатады. Бір жағынан, ұлттық стратегиялар, жасыл таксономияның
жекелеген элементтері, AIFC Green Finance Centre қызметі мен даму
институттарының қатысуы ESG қағидаттары мен «жасыл» экономиканың тұрақты
даму саясатына кезең-кезеңімен ендіріліп жатқанын көрсетеді {[}6, 11{]}.
Екінші жағынан, верификация рәсімдерінің бытыраңқылығы, ESG-есептіліктің
бірыңғай стандарттарының кеңінен қолданылмауы, әсер көрсеткіштерін
айқындау әдістемелерінің бірыңғайланбауы гринвошинг тәуекелдерін
күшейтіп, инвесторлардың сеніміне теріс ықпал етуі ықтимал {[}5, 9{]}.
Осы тұрғыдан алғанда, халықаралық әдебиетте талқыланатын корпоративтік
ESG-тәжірибенің шынайылығы мен «маркетингтік» сипаттағы дискурс
арасындағы қайшылықтар Қазақстан контексті үшін де өзекті болып отыр.

Гринвошинг тәуекелдері мен оларды шектеу жолдарына қатысты алынған
нәтижелер халықаралық талқылауларда айқындалған негізгі бағдарлармен
сәйкес келеді. Әдебиеттерде ESG-инвестициялардың көлемі жедел ұлғайған
сайын компаниялар мен қаржы ұйымдары ESG-тілдес риториканы кеңінен
қолдана бастайтыны, алайда қатаң реттеу мен біріздендірілген есептілік
болмаған жағдайда бұл үрдіс формализмге ұласу қаупін күшейтетіні
көрсетіледі {[}9{]}. Қазақстан жағдайында жасыл облигациялар мен
несиелер бойынша жобалардың «жасыл» мәртебесін айқындау кезінде
анықталған олқылықтар сыртқы верификаторлар институтын, тәуелсіз
ESG-рейтингтерді және ашық есептілік талаптарын күшейту қажеттілігін
айғақтайды. Бұл, өз кезегінде, ICMA Green Bond Principles, Climate Bonds
Standard және өзге де халықаралық стандарттарға жақындасу мен оларды
ұлттық деңгейде бейімдеу арасындағы оңтайлы тепе-теңдікті табуды талап
етеді {[}12{]}.

ESG қағидаттарын ұлттық тұрақты даму стратегияларымен сәйкестендіру
мәселесіне қатысты жүргізілген талдау Қазақстанның бұл бағыттағы
мүмкіндіктері айтарлықтай екенін көрсетті. Соңғы жылдары жарияланған
еңбектерде ESG-саясатты ұлттық деңгейде енгізу инвестициялық
тартымдылықты күшейтуге, институционалдық ортаның сапасын арттыруға және
қоғамның жалпы әл-ауқатын жоғарылатуға үлес қосатыны дәлелденген
{[}11{]}. Біздің зерттеу қорытындылары жасыл қаржыландыру құралдарын
(жасыл облигациялар, жасыл несиелер, әлеуметтік және ESG-облигациялар)
ұлттық стратегиялық құжаттармен, декарбонизация міндеттерімен және
«жасыл» күн тәртібімен нақты түрде үйлестіру қажеттігін көрсетіп отыр.
Алайда бұл өзара байланыс әлі толық қалыптасқан жоқ: стратегиялық
құжаттар көбіне жалпы бағдар берумен шектелсе, қаржылық өнімдер мен
реттеуші механизмдер көбінесе пилоттық немесе жекелеген жобалар
деңгейінде іске асырылуда.

Экономикалық және әлеуметтік тиімділік призмасы арқылы жасыл
облигациялар мен несиелердің әсерін талдау халықаралық әдебиетте
жасалған тұжырымдармен бірқатар ортақ белгілерге ие екенін көрсетті.
Эмпирикалық зерттеулерде green bond issuance компаниялар үшін капитал
құнының төмендеуіне, ақпараттық ашықтықтың артуына және нарықтағы
беделдің нығаюына алып келетіні дәлелденген {[}8{]}. Біздің талдау да
эмитенттер үшін ұзақ мерзімді қаржыландыру базасын кеңейту,
ESG-тәжірибелерді ұйымдық деңгейде бекіту және мүдделі тараптар
алдындағы есептілік дәрежесін жоғарылату сияқты оң нәтижелерді
айқындады. Инвесторлар жағынан қарағанда, тәуекелдерді әртараптандыру,
ұзақ мерзімді тұрақтылық және қаржылық, сондай-ақ әлеуметтік-экологиялық
әсерді біріктіретін «қосарланған табыстылық» тұжырымдамасы негізгі
тартымдылық факторларының бірі болып отыр.

Сонымен бірге, алынған нәтижелер зерттеудің шектеулері мен алдағы
ізденістерге қатысты бірқатар бағыттарды да айқындады. Біріншіден,
деректердің қолжетімділігі мен толықтығы шектеулі: жасыл және
ESG-облигациялар, жасыл несиелер, сондай-ақ жоба деңгейіндегі әсер
индикаторлары бойынша ашық статистика жеткілікті түрде жинақталмаған.
Екіншіден, осы жұмыста қолданылған тәсіл негізінен сипаттамалық және
салыстырмалы талдаумен шектеледі; келешекте жасыл қаржыландырудың
эмитенттердің қаржылық нәтижелеріне, капитал құнына және көміртек ізі
динамикасына ықпалын эконометрикалық модельдер арқылы сандық бағалау
маңызды болып табылады. Үшіншіден, инвесторлардың жарияланған
ESG-ұстанымдары мен олардың нақты инвестициялық таңдаулары арасындағы
алшақтықты, сондай-ақ халықтың «жасыл» қаржы құралдары бойынша қаржылық
сауаттылық деңгейін арнайы зерттеу қажет.

Жалпы, алынған деректер зерттеудің бастапқы гипотезасын белгілі бір
дәрежеде қуаттады: егер Қазақстанда ESG қағидаттарына негізделген жасыл
облигациялар мен жасыл несиелер нарығының институционалдық және
реттеушілік іргетасы нығайып, верификация мен ашықтық механизмдері
дамытылса, бұл құралдар «жасыл» экономикаға көшу үдерісін жеделдетіп,
ұлттық тұрақты даму мақсаттарына жетуге мәнді үлес қоса алады. Сонымен
қатар, осы трансформацияның нәтижелілігі жасыл қаржыландыруды жалпы
экономикалық саясатпен, салалық даму стратегияларымен және қоғамда
жауапты инвестициялау мәдениетін қалыптастыру деңгейімен қаншалықты
үйлесімді ұштасатынына тәуелді болады.

{\bfseries Қорытынды.} Зерттеудің негізгі мақсаты -Қазақстан контексінде
ESG қағидаттарына сүйенетін жасыл облигациялар, жасыл несиелер және
инвестициялар институттарының теориялық-әдіснамалық алғышарттарын
айқындау, олардың жасыл экономиканы қалыптастырудағы және тұрақты дамуды
қамтамасыз етудегі маңызын талдау, сондай-ақ бұл құралдарды кейінгі
кезеңде ілгерілетуге бағытталған ұсынымдар ұсыну болды. Қойылған
мақсатқа жету үшін жүйелік, институционалдық және салыстырмалы талдау
әдіснамасы қолданылып, ғылыми еңбектерге шолу жасалды,
нормативтік-құқықтық және стратегиялық құжаттарға контент-талдау
жүргізілді, халықаралық және ұлттық статистикалық мәліметтер,
эмитенттердің ашық ESG-есептілігі мен жасыл облигациялар эмиссиясына
қатысты ақпарат талданды.

Алынған деректер бірқатар маңызды ғылыми тұжырым жасауға негіз болды.
Біріншіден, ESG өлшемдері мен қаржылық сипатын ескеріп құрылған жасыл
қаржылық құралдардың матрицалық жіктемесі жасыл облигациялар, несиелер,
әлеуметтік және ESG-облигациялар, жасыл инвестициялық қорлар мен
микрокредиттердің өзара байланысы мен орнына түсінік беріп, ұлттық жасыл
таксономияны қалыптастыруға әдістемелік база ұсынды. Екіншіден, әлемдік
және қазақстандық жасыл/ESG-облигациялар нарықтарының динамикасын
салыстыру нәтижесінде жаһандық деңгейде бұл сегменттің тұрақты өсу
бағытын сақтап отырғаны, ал Қазақстанда нарықтың кеш пайда болғанымен,
соңғы жылдары өте жоғары өсім қарқынына жеткені анықталды; сонымен
қатар, жасыл құралдардың корпоративтік облигациялар нарығындағы үлесі
әлі де аз екені байқалды. Үшіншіден, жаңартылатын энергетика, энергия
тиімділігі, коммуналдық инфрақұрылым, көлік, су ресурстары және
агросектор жасыл қаржыландырудың негізгі басым салалары ретінде
белгіленіп, МЖӘ жобаларының қаржыландыру құрылымында жасыл компоненттің
әзірге қосалқы, бірақ әлеуеті жоғары бағыт екені көрсетілді.
Төртіншіден, нормативтік базаның біртіндеп дамуына қарамастан,
верификация, ESG-есептілік стандарттары және әсер көрсеткіштерін есептеу
әдістемелеріндегі кемшіліктер гринвошинг тәуекелдерін күшейтетіні
анықталып, осыған байланысты институционалдық тетіктерді жетілдіру
қажеттігі негізделді.

Осы қорытындылар бастапқы гипотезаның жалпы алғанда расталатынын
көрсетеді: Қазақстанда ESG қағидаттарына негізделген жасыл облигациялар
мен жасыл несиелер нарығының институционалдық және реттеушілік базасы
нығайып, верификация және ашықтық тетіктері дамыған жағдайда, аталған
құралдар «жасыл» экономикаға көшу үдерісін жеделдетуге, ұлттық тұрақты
даму мақсаттарына қол жеткізуге, сондай-ақ қаржы жүйесінің орнықтылығын
күшейтуге айтарлықтай үлес қоса алады. Бұдан бөлек, зерттеу жасыл
қаржыландырудың экономикалық (капитал құнын азайту, инвесторлар шоғырын
кеңейту, беделдік тиімділік), экологиялық (көміртек ізін төмендету,
ресурстарды ұтымды пайдалану) және әлеуметтік (жұмыс орындарын құру,
инфрақұрылым сапасын жақсарту) нәтижелерін кешенді түрде бағалау және
оларды саясаттық әрі корпоративтік шешімдерге жүйелі түрде енгізу
қажеттігін айқындап отыр.

Перспективаларды бағалау тұрғысынан, біріншіден, ұлттық жасыл
таксономияны халықаралық стандарттармен сәйкестендіре отырып толықтыру,
жобалардың «жасыл» мәртебесін тәуелсіз верификаторлар арқылы растау
тетігін қалыптастыру, ESG-есептіліктің бірыңғай үлгілерін енгізу және
жасыл облигациялар мен несиелер бойынша деректерді қамтитын ашық
статистикалық база жасау қажеттілігі алға шығады. Екіншіден, болашақ
зерттеулерде жасыл қаржыландырудың эмитенттердің қаржылық нәтижелеріне,
капитал құнына және көміртек шығарындыларына ықпалын эконометрикалық
тәсілдермен өлшеу, сондай-ақ инвесторлардың жарияланған ESG-ұстанымдары
мен олардың нақты инвестициялық таңдаулары арасындағы алшақтықты
эмпирикалық талдау өзекті. Үшіншіден, жасыл қаржы құралдарын МЖӘ
механизмдерімен, мемлекеттік бюджет және даму институттары
бағдарламаларымен тығыз ұштастыру, өңірлік және жергілікті деңгейде
«жасыл» жобаларға арналған пилоттық платформалар құру, халық пен бизнес
үшін қаржылық әрі экологиялық сауаттылықты арттыруға бағытталған кешенді
шаралар пакетін әзірлеу маңызды. Осындай шаралар жасыл қаржыландыру
инфрақұрылымын Қазақстанның экономикалық, экологиялық және әлеуметтік
даму стратегияларымен терең интеграциялауға және алынған ғылыми
нәтижелерді нақты саясат пен тәжірибеде іске асыруға жағдай жасайды.
\end{multicols}

\begin{center}
{\bfseries Әдебиеттер}
\end{center}

\begin{refs}
1. International Capital Market Association (ICMA). Green Bond
Principles (GBP).2025 edition.
-URL:https://www.icmagroup.org/sustainable-finance/the-principles-guidelines-and-handbooks/green-bond-principles-gbp/-
Қаралған күні: 15.01.2026.

2. AIFC Green Finance Centre. Green finance market in Kazakhstan.
Analytical report. Astana: AIFC Authority. -2023.-URL
https://aifc.kz/news/green-finance-market-in-kazakhstan/- Қаралған күні:
15.01.2026.

3. United Nations Development Programme in Kazakhstan. UNDP catalyses
issuance of first green bonds in Kazakhstan.29.10.2020.-URL:
https://www.undp.org/kazakhstan/stories/undp-catalyses-issuance-first-green-bonds-kazakhstan-
Қаралған күні: 15.01.2026

4. Петросянц Т.В. Современное состояние «зеленого» финансирования в
Казахстане // Colloquium-journal. -2023. -№ 29(188). -С.32-33.DOI
10.24412/2520-6990-2023-29188-32-33.

5. Мангибаева Д.Д. Потенциал развития инструментов зеленого
финансирования в Республике Казахстан // Colloquium-journal. -2023.
-№28(173). -С.39-43. DOI 10.24412/2520-6990-2023-28187-39-43

6. Жагыпарова А.О., Мухияева Д.М., Сембиева Л.М. Қазақстан үшін «жасыл»
экономиканың басымдықтары // Education. Quality Assurance. -2025. -№ 1.
-Б.93-103. DOI 10.58319/26170493\_2025\_1\_93/

7. Калмыков Д.Е. Жасыл экономика бойынша қысқаша нұсқаулық.
ЭкоМұражай//Sinest. -2017.-32 б.

8. Zhang R., Li Y., Liu Y. Green bond issuance and corporate cost of
capital//Pacific-Basin Finance Journal. -2021.-Vol.69:01626. DOI
10.1016/j.pacfin.2021.101626.

9. Kräussl R., Oladiran T., Stefanova D. A review on ESG investing:
Investors' expectations, beliefs and perceptions//Journal of Economic
Surveys. -2024.-Vol.38(2). -P.476-502. DOI 10.1111/joes.12599.

10. Yessymkhanova Z., Birmagambetov T., Dauletkhanova Zh., Baimamyrov B.
Current status of green finance in Kazakhstan // E3S Web of Conferences.
-2023.-Vol.402:08033. DOI 10.1051/e3sconf/202340208033.

11. Spankulova L., Aben A., Nurmukhametov N., Zeinolla S. Integration of
ESG principles into the sustainable development strategy of Kazakhstan
// Journal of Central Asian Studies. -2025.-Vol.23(1) - P.61-86. DOI
10.52536/3006-807X.2025-1.004.

12. Сlimate bonds standard Globally recognised, Paris-aligned
Certification of debt instruments, entities and assets using robust,
science-based methodologies.-
\url{URL:https://www.climatebonds.net/files/documents/Climate-Bonds_Climate-Bonds-Standard_V4-1_Feb-2024.pdf}
. - Қаралған күні: 15.01.2026.
\end{refs}

\begin{center}
{\bfseries References}
\end{center}

\begin{refs}
1. International Capital Market Association (ICMA). Green Bond
Principles (GBP).2025 edition.
-URL:https://www.icmagroup.org/sustainable-finance/the-principles-guidelines-and-handbooks/green-bond-principles-gbp/
- Date of access: 15.01.2026.

2. AIFC Green Finance Centre. Green finance market in Kazakhstan.
Analytical report. Astana: AIFC Authority. -2023.-URL
https://aifc.kz/news/green-finance-market-in-kazakhstan - Date of
access: 15.01.2026.

3. United Nations Development Programme in Kazakhstan. UNDP catalyses
issuance of first green bonds in Kazakhstan.29.10.2020.-URL:
\url{https://www.undp.org/kazakhstan/stories/undp-catalyses-issuance-first-green-bonds-kazakhstan}
- Date of access: 15.01.2026.

4. Petrosjanc T.V. Sovremennoe sostojanie «zelenogo» finansirovanija v
Kazahstane // Colloquium-journal. -2023. -№ 29(188). -S.32-33.DOI
10.24412/2520-6990-2023-29188-32-33. {[}in Russian{]}

5. Mangibaeva D.D. Potencial razvitija instrumentov zelenogo
finansirovanija v Respublike Kazahstan // Colloquium-journal. -2023.
-№28(173). -S.39-43.DOI 10.2441/2526-6990-2023-28187-39-43. {[}in
Russian{]}

6. Zhagyparova A.O., Muhijaeva D.M., Sembieva L.M. Қazaқstan үshіn
«zhasyl» jekonomikanyң basymdyқtary // Education. Quality Assurance.
-2025. -№ 1. -B.93-103. DOI 10.58319/26170493\_2025\_1\_93/ {[}in
Russian{]}

7. Kalmykov D.E. Zhasyl jekonomika bojynsha қysқasha nұsқaulyқ.
JekoMұrazhaj//Sinest. -2017.-32 b. {[}in Kazakh{]}

8. Zhang R., Li Y., Liu Y. Green bond issuance and corporate cost of
capital//Pacific-Basin Finance Journal. -2021.-Vol.69:01626. DOI
10.1016/j.pacfin.2021.101626.

9. Kräussl R., Oladiran T., Stefanova D. A review on ESG investing:
Investors' expectations, beliefs and perceptions//Journal of Economic
Surveys. -2024.-Vol.38(2). -P.476-502. DOI 10.1111/joes.12599.

10. Yessymkhanova Z., Birmagambetov T., Dauletkhanova Zh., Baimamyrov B.
Current status of green finance in Kazakhstan // E3S Web of Conferences.
-2023.-Vol.402:08033. DOI 10.1051/e3sconf/202340208033.

11. Spankulova L., Aben A., Nurmukhametov N., Zeinolla S. Integration of
ESG principles into the sustainable development strategy of Kazakhstan
// Journal of Central Asian Studies. -2025.-Vol.23(1) - P.61-86. DOI
10.52536/3006-807X.2025-1.004.

12. Сlimate bonds standard Globally recognised, Paris-aligned
Certification of debt instruments, entities and assets using robust,
science-based methodologies.-
\url{URL:https://www.climatebonds.net/files/documents/Climate-Bonds_Climate-Bonds-Standard_V4-1_Feb-2024.pdf}
. - Date of access: 15.01.2026.
\end{refs}

\begin{info}
\hspace{1em}\emph{{\bfseries Авторлар туралы мәліметтер}}

Мамутова Қ.Қ.-экономика ғылымдарының кандидаты, М.Әуезов атындағы
Оңтүстік Қазақстан университеті, Шымкент, Қазақстан, e-mail:
Katira.7575@mail.ru;

Машраб Ж.С.- докторанты, М.Әуезов атындағы Оңтүстік Қазақстан
университеті, Шымкент, Қазақстан,e-mail: zmashrab@bk.ru;

Айтымбетова А.Н.-экономика ғылымдарының кандидаты, қауымдастырылған
профессор, М.Әуезов атындағы Оңтүстік Қазақстан университеті, Шымкент,
Қазақстан, e-mail: a.ainura-81@mail.ru.

\hspace{1em}\emph{{\bfseries Information about the authors}}

Mamutova K.- Candidate of Economic Sciences, associate professor,
Auezov University, Shymkent, Kazakhstan, е-mail: Katira.7575@mail.ru;

Mashrab Zh.- PhD student, Auezov University, Shymkent, Kazakhstan,
е-mail: zmashrab@bk.ru;

Aitymbetova A.- Candidate of Economic Sciences associate professor,
Auezov University, Shymkent, Kazakhstan, е-mail: a.ainura-81@mail.ru.
\end{info}
